\section{Bundles on a tangent box}\label{sec:transport}

A tangent box separates the two bundles as well as the two tangent
directions. We show that each pulled-back summand comes from its
corresponding factor. Transport from the diagonal gives a natural
normalization of this identification at the point where the bracket
will be tested.

For \(i\ne j\), the splitting defines a projected-bracket partial
connection on \(E_i\) in directions of \(E_j\):
\begin{equation}\label{eq:partial}
\nabla^i_\xi e=P_i[\xi,e],
\end{equation}
where \(\xi\) and \(e\) are local holomorphic sections of \(E_j\)
and \(E_i\), respectively.

\Needspace{7\baselineskip}
\begin{proposition}\label[proposition]{prop:transport}
Let \(F:S\to X\) be a tangent box of \(g:C\to X\). Transport along
the fibers of \(\pr_i\) gives holomorphic bundle isomorphisms
\begin{equation}\label{eq:transport}
\Theta_i:\pr_i^*(g^*E_i)\xrightarrow{\ \sim\ }F^*E_i,
\qquad i=1,2,
\end{equation}
whose restrictions to the diagonal are the identity.
\end{proposition}

\begin{proof}
The bracket rules and \(P_i\xi=0\) give
\[
\nabla^i_{f\xi}e=f\nabla^i_\xi e,\qquad
\nabla^i_\xi(fe)=\xi(f)e+f\nabla^i_\xi e.
\]
Thus \eqref{eq:partial} is tensorial in its \(E_j\) direction and obeys
the connection Leibniz rule in its \(E_i\) argument. These two rules
use no bracket of two \(E_j\) sections and require no integrability
of either summand.

Fix \(i\) and let \(j\ne i\). The tangency of \(F\) sends the vertical
line bundle \(\ker(d\pr_i)=T^{(j)}\) into \(F^*E_j\). The partial
connection therefore pulls back to a holomorphic connection on
\(F^*E_i\) relative to \(\pr_i\). We spell out this construction because
\(F\) may have critical points. Take local frames \(e_\lambda\) of
\(E_i\) and \(\xi_\gamma\) of \(E_j\), and write
\[
\nabla^i_{\xi_\gamma}e_\lambda
 =\sum_\nu\Gamma_{\gamma\lambda}^{\nu}e_\nu,\qquad
dF(V)=\sum_\gamma h_\gamma F^*\xi_\gamma
\]
for a local holomorphic vertical vector field \(V\) on \(S\). Define
\[
D_V\!\left(\sum_\lambda f_\lambda F^*e_\lambda\right)
=\sum_\lambda V(f_\lambda)F^*e_\lambda
 +\sum_{\gamma,\lambda,\nu}
 f_\lambda h_\gamma
 (\Gamma_{\gamma\lambda}^{\nu}\circ F)F^*e_\nu.
\]
The two connection rules and the chain rule
\(V(A\circ F)=dA(dF(V))\) give the usual connection transformation
under a change of the \(E_i\) frame; tensoriality in the first argument
handles a change of the \(E_j\) frame. Hence these formulas agree on
overlapping product neighborhoods. They define a holomorphic relative
connection
\[
D:\mathcal O(F^*E_i)\longrightarrow
\mathcal O(F^*E_i\otimes\Omega^1_{S/C}),
\]
where the relative differential is for \(\pr_i:S\to C\). At a point where
\(dF(V)=0\), all \(h_\gamma\) vanish there and the formula still applies;
no inverse of \(dF\) is used.

The restriction of \(D\) to any fiber \(C=\PP^1\) is a holomorphic
connection. Its curvature is a holomorphic two-form on that curve and
therefore vanishes. Horizontal sections satisfy an everywhere nonsingular
holomorphic linear ordinary differential equation. Local existence and
uniqueness give continuation along every path, with reverse transport as
its inverse. Flatness and simple connectedness of \(\PP^1\) make this
transport path-independent. Starting at the fiber's diagonal point
therefore identifies its bundle with the corresponding fiber of
\(g^*E_i\).

These identifications depend holomorphically on both variables. Fix a
base point and a path in its fiber from the diagonal to a desired endpoint.
Cover the path by a finite chain of coordinate disks and bundle frames,
and shrink the base neighborhood so that the same chain applies to nearby
fibers. In each disk the coefficient matrix of the horizontal equation
is holomorphic in the base and fiber coordinates. Picard iteration
converges normally on smaller disks by the exponential-series estimate,
so its solution matrix is holomorphic in the base and the start and end
coordinates. Composing these local transports and frame changes gives a
holomorphic identification near the chosen point of \(S\). Path
independence makes the local identifications agree globally. Transport
at the starting diagonal point is the identity, giving \(\Theta_i\)
and its stated normalization.
\end{proof}
